% =============================================================================
% Project : THE_PHILOSOPHY_OF_VIABILITY
% Book    : Book I — Ontology of Viability
% Book ID : VS.B01
%
% File    : 4_ASSETS/tables/
%           VS.B01.P01.C03.S05.SS03.P012.TL01.components_composition_whole_and_identity.tex
%
% Asset   : TL01 — Components, Composition, Whole, and Identity
%
% Purpose :
%   To distinguish the principal ontological terms involved
%   in compositional analysis and to prevent unsupported
%   inferences from component multiplicity to identity.
%
% =============================================================================

\begin{table}[htbp]
\centering
\caption{Components, Composition, Whole, and Identity}
\label{tab:components-composition-whole-and-identity}

\small
\renewcommand{\arraystretch}{1.35}

\begin{tabular}{p{0.19\textwidth} p{0.34\textwidth} p{0.35\textwidth}}
\hline
\textbf{Term}
&
\textbf{Primary concern}
&
\textbf{What it does not establish by itself}
\\
\hline

Component
&
An element participating in a broader configuration.
&
That it possesses an independently distinct identity.
\\

Composition
&
The organized relations through which components
participate in a configuration.
&
That the resulting configuration possesses exactly
one identity.
\\

Aggregate
&
A plurality of elements considered together under
a grouping or unity criterion.
&
That the grouping constitutes a composition or
a distinct identity.
\\

Composed whole
&
A configuration considered through the organization
of its participating components.
&
That its identity is reducible to its components
or their arrangement.
\\

Structural organization
&
The positions and relations through which components
are organized.
&
That structural unity is equivalent to identity unity.
\\

Material constitution
&
The material elements and conditions involved
in the constitution of an existent.
&
That material continuity or integration determines
identity.
\\

Identity
&
The ontological unity or distinction attributable
at the level of identity.
&
That its status can be inferred solely from
component count or compositional organization.
\\

\hline
\end{tabular}

\end{table}

% =============================================================================
% END OF ASSET TL01
% =============================================================================